\chapter{Discussion and Conclusion}\label{ch:discussion}

\section{Where the Architecture Binds}\label{sec:disc-rq1}

\textsc{argued}, from Section~\ref{sec:res-judge}.
MMLU-Pro is multiple-choice with a verifiable answer key. The judge is
shown the options and, in most cases, a response that states its own
selected option (Section~\ref{sec:arch-judge}). A judge capable of
solving the item can therefore verify both responses against its own
solution and decide the comparison without consulting the rubric at
all --- while still being answer-key-blind in the sense the protocol
guarantees, because nothing in the prompt tells it which option is
right.

The evidence that this is what it does is
Section~\ref{sec:res-judge-summary}: five measurements, four with
matched controls, all pointing the same way. None of the numbers is
restated here; this section draws only the inference.

The consequence for the architecture is precise and it is not a
failure. \textbf{The partition can only bind where correctness is not
checkable.} On open-ended generation, on agent traces, on tasks with no
key --- the settings the routing motivation of
Section~\ref{sec:motivation} is actually about --- the judge has no
shortcut available and the rubric is the only thing standing between it
and an undocumented heuristic. On a multiple-choice corpus it has a
shortcut, and it takes it.

The two halves of that statement belong together and must be stated
together: \emph{MMLU-Pro
grounds the comparison against ground truth, and is simultaneously the
corpus on which the judging mechanism can bypass what the comparison is
about.} The first is why the arena could be validated at all. The
second is why validating it there does not establish the thing the
architecture was built for.

This is an inference from the Math and law runs, not a demonstration.
The demonstration is the options-blind arm --- the same questions, the
same rubric, the options withheld --- whose predictions are already
fixed in writing and which is blocked on a trace-complete roster
(Section~\ref{sec:disc-future}).

\section{What Each Negative Means, Link by Link}\label{sec:disc-rq2}

\textsc{argued}. This thesis reports several negative results, and they
do not mean the same thing, because they land on different links of the
chain in Section~\ref{sec:research-questions}.

\textbf{``Cells do not carry different true rankings'' lands on link 3's
premise.} It is a statement about MMLU-Pro and about a roster spanning
$54$ accuracy points: cell identity adds essentially nothing to model
ordering beyond global capability. It says nothing at all about whether
cell-scoped judging produces better judgments, because it never touches
a verdict --- it is computed on ground-truth accuracy alone.

\textbf{``The rubric does not change the verdict'' lands on link 3's
mechanism.} And it is not a parity finding. Two judges that agree on
$99.4\%$ of winners cannot produce a detectable ranking difference at
any sample size, granularity or domain count. The arms barely differ,
and no amount of power repairs a comparison whose treatment is nearly
the control. The transferable consequence is a precondition rather than
a result: measure verdict agreement before spending a judging budget on
an A/B arm.

\textbf{Neither touches links 1 or 2, and both of those came out
positive.} The construction reaches a certified fixed point and
produces cells of the intended size distribution; rubrics induced on
those cells are specific to them against a size-matched randomised null
at $p = 1.21 \times 10^{-6}$, with all four pre-registered
discriminators firing in the predicted directions. The chain does not
break at the partition. It breaks at the step where a better rubric was
supposed to produce a better verdict, on a corpus where the verdict was
already decided by something else.

\textbf{And the reading rule matters more than the reading.} A null
under a decision rule fixed before the data is a finding. A null under a
rule chosen afterwards is a shrug. Every negative above was read
against a criterion recorded in a dated document: three outcomes named
in advance for the granularity analysis, with ``observed $\rho$ is
primary'' fixed in advance --- and that mattered, because the
disattenuated column moves the opposite way and clips past $1.0$, so
leading with it would have produced the opposite conclusion from the
arithmetic of the correction rather than from the data. Four
directional predictions with the confound direction named per measure
for the rubric null. A decisive difference of two Spearman grid steps,
and four void conditions, fixed for the no-partition baseline before its
verdicts existed.

A reader must not be allowed to collapse ``cells do not carry different
true rankings'' into ``the taxonomy does not help the arena''. The first
is a property of this corpus's saturation for this roster. The second is
a claim about judge quality, and the measurement that would decide it
has not been run in the only regime where the rubric has been shown able
to matter.

\section{The Apparatus Caught Its Own Errors}\label{sec:disc-selfcheck}

\textsc{measured}, four instances. None of the following was found by
looking for it, and that is the point: they are the best available
evidence that the results the apparatus did \emph{not} overturn are
sound.

\textbf{The response-capture gap} was found while sizing the roster for
a different experiment. It showed that half of the pilot arena's
comparisons had been decided by which side had text at all, and it
qualified figures that were already written down
(Section~\ref{sec:res-capture-gap}).

\textbf{A side-channel in two candidate roster entries} was found by a
pre-launch format check whose only anomalous column was newline count:
the two systems' stored output was a raw JSON envelope carrying a
self-reported code that predicts the very outcome being judged.
Neither system entered any arena run as stored; the full account is in
Appendix~\ref{app:corrections-register}.

\textbf{The reliability constant} was caught twice, by two independent
routes that compounded to the same factor: a re-fit at a different
roster size, and a units check on a split-half length mismatch
(Section~\ref{sec:disc-debt}).

\textbf{The domain screen's two reversals} both traced to one cause,
found only because the screen was re-run at the roster the arena would
actually use (Section~\ref{sec:res-law}).

Two design features did the work in three of these four cases, and both
are cheap. A check whose only informative column is an unglamorous one
--- newline count --- caught what length, citation rate and URL rate
did not. And re-deriving a statistic at the configuration the decision
will actually run at, rather than at whatever configuration was
convenient, caught both screen reversals.

\section{Threats to Validity and Limitations}\label{sec:disc-validity}

\subsection{Internal Validity}\label{subsec:disc-internal-validity}

The information separation that the answer-key-blind claim rests on is
enforced in code and audited by a build-failing assertion
(Section~\ref{sec:arch-judge}), which is the strongest guarantee in the
system. It does not, however, make the judge answer-\emph{blind}, and
Section~\ref{sec:disc-rq1} is the consequence.

Two internal threats are unmitigated. Comparisons are not independent
given the Bradley--Terry parameters, because a single held-out query is
reused across many pairs; the Fisher standard errors treat them as
independent and therefore understate uncertainty by an unmeasured
margin. And the scheduler over-samples uncertain pairs, which biases
win-rate (Section~\ref{sec:meth-scheduling}); $\theta$ is quoted
throughout.

\subsection{External Validity}\label{subsec:disc-external-validity}

Every arena result in this thesis is on one corpus, at seed 42, on two
domains (Math and law; replications on philosophy and history are
launched with results pending), with rosters of 8 to 12 models. The construction is
contingent on one embedding model: if the embedding changes, the
snapshot is invalid and results produced under different embeddings are
not comparable. The anchor structure is seeded from MMLU-Pro's
editorial domains and thereafter immutable, so anchor-level results are
rankings conditioned on MMLU-Pro's partitioning; cross-domain migration
mitigates this at the query level but cannot relocate or merge anchors.
Discovery is bounded by the corpus available at construction time,
and there is no incremental insertion mechanism --- a new corpus
requires a full reconstruction.

The roster exclusion policy is a further external-validity limit. Ten
models are excluded by a hard-coded list of names, each with a
recorded reason; the policy is a ratchet on known failures, not a
threshold, and it does not generalise to any new model
(Section~\ref{sec:exp-rosters}). A general guard is a one-line change
and is named as future work.

\subsection{Construct Validity}\label{subsec:disc-construct-validity}

\textbf{Judge-side benchmark contamination.}
MMLU-Pro is a public benchmark, and its questions and answer keys
plausibly appear in the pre-training corpora of contemporary
LLMs~\parencite{deng2024contamination}; contamination-limited benchmark
designs exist precisely because the problem is
pervasive~\parencite{white2024livebench}.
The specific threat here is not that the \emph{evaluated} models saw
the questions --- that affects the accuracy oracle and the arena
symmetrically --- but that the \emph{judge} may recognise a held-out
question and its answer from training. A memorising judge approximates
the accuracy oracle directly, and would inflate rank agreement for
reasons unrelated to the evaluation mechanism. The answer-key-blind
protocol does not defend against it, because the information sits in
the judge's weights rather than in the prompt.

The generic-judge and no-partition arms discipline this threat only
partially, and less than an earlier version of this argument claimed.
Because all arms share the same judge model, contamination contributes
to each alike and is differenced out of any contrast between them --- so
a rubric gain that survived the control could not be explained by
memorised answer keys. But the measured contrast is approximately zero
(Sections~\ref{sec:res-generic-judge} and~\ref{sec:res-c5}), so there is
no gain for the control to protect. What the control establishes is
therefore the absence of an effect, not the contamination-robustness of
a present one. The absolute agreement levels remain an upper bound on
mechanism-attributable validity, and a direct estimate of judge
memorisation --- verdict consistency under paraphrase, or a probe set
published after the judge's training cutoff --- is future work.

\textbf{Residual judge bias.}
The dual-call protocol converts positional disagreement into an
explicit tie, and a post-hoc adjustment rewrites a pair's counts when
its order asymmetry is large (Section~\ref{sec:arch-judge}). Verbosity
bias~\parencite{chen2024verbosity} is not mitigated: the rubric
contains no length-normalisation criterion, and the capture-gap result
of Section~\ref{sec:res-capture-gap} is an extreme instance of a
length-adjacent preference --- on mixed pairs the judge takes the side
with text on $863/868$ of comparisons. Self-preference
bias~\parencite{panickssery2024llm} does not arise here, because the
judge model appears in no roster,\footnote{Verified rather than
assumed: the judge is Mistral-Large-3 in every reported run, confirmed
against the Run~B run log (four occurrences of the judge identifier,
none of any smaller ``ministral'' variant) and the runtime
configuration; no roster contains it (Appendix~\ref{app:models}).}
but no constraint enforces that; it
is a property of these runs rather than of the system.

\textbf{Tie approximation.}
Ties are encoded as half-wins inside the standard Bradley--Terry
likelihood, without Davidson's tie-propensity
parameter~\parencite{davidson1970ties}. Section~\ref{sec:res-tie-policy}
shows this is not a harmless approximation: on one paired contrast the
choice of tie convention decides which arm appears ahead. A Davidson
fit would make the tie propensity a parameter rather than a convention
and is the principled repair.

\textbf{Budget accounting.}
Comparison counts follow the convention of Chapter~\ref{ch:results}:
$B_{\text{logical}}$, with $B_{\text{API}} = 2 B_{\text{logical}}$
under the dual-call design, so that cost figures remain comparable
with single-call systems.

\subsection{The Standing Measurement Debt}\label{sec:disc-debt}

Two of the following are errors this thesis found in its own
previously reported numbers. They are presented under the frame of
Section~\ref{sec:disc-selfcheck} rather than apologetically. The full
account of each --- the derivation, the pre-correction values, and
what was withdrawn --- is told once, in the corrections register of
Appendix~\ref{app:corrections-register}; no interval computed before
the corrections is quoted anywhere in this document. What remains here
is the corrected state and where each correction binds.

\begin{itemize}

\item \textbf{Bradley--Terry standard errors.} The rank completion
requires subtracting $1/K^2$ after inversion; the implementation
subtracted $1/K$, so every pre-correction standard error was the floor
constant, and correct values are 44 to 209 times larger
(Section~\ref{sec:meth-bt}). No pre-correction interval or ``resolved
pair'' determination is quoted.

\item \textbf{The scheduler consequence of that bug, and which runs it
touches.} The fix is dated (commit \texttt{a92a3ed}, 2026-07-27), and
Run~A predates it. Run~A's scheduler therefore consumed floor-constant
standard errors, with three mechanical consequences: the variance
division inside $U_{\text{resolve}}$ degenerated to a constant, so the
term collapsed to pure entropy; the maturity-driven mixing weight
$\alpha$, keyed to falling per-leaf standard errors, never matured;
and the $3\sigma$ resolve rule reduced to a fixed margin
$|\hat{s}_i - \hat{s}_j| \geq 0.06$. Run~A's comparison allocation was
consequently near-uniform-entropy scheduling, not the
variance-weighted design described in
Section~\ref{sec:meth-scheduling}, and no reading of Run~A may lean on
adaptive allocation. Runs~B and~C postdate the fix and ran with
measured standard errors.

\item \textbf{The reliability constant and the disattenuation form.}
The published $c = 7.66$ corrects to $c = 2.52$ --- a split-half
length mismatch compounded with a roster re-fit --- and disattenuation
is two-sided, $\rho / r$ rather than $\rho / \sqrt{r}$
(Section~\ref{sec:exp-metrics}). The correction compresses the
reliability ladder across granularities from a spread of $0.219$ to
$0.088$, so leaf-scale cells stop being individually low-reliability;
that weakens the reliability leg of the argument for judging at a
coarser cut, and the identification-cost result of
Section~\ref{sec:res-cost} is untouched and probably carries that
argument alone.

\item \textbf{The evaluative-redundancy percentages} could not be
reproduced under either an observed-primary or a disattenuated
reading and are withdrawn; no redundancy percentage is asserted in
this thesis.

\item \textbf{Two convention-dependent figures} are reported as such
rather than quoted flat: the correctness-blind rank correlation
(Section~\ref{sec:res-correctness}) and the leaf-level between-cell
agreement, $0.922$ against $0.884$ under two defensible assignment
conventions (Section~\ref{sec:res-granularity}); the governing
convention is named where each number appears.

\item \textbf{Routing expected calibration error} is exported under a
mis-specified aggregation --- a maximum over leaves where a sum is
required, and a maximum over leaves is not a distribution. The
exported $0.2114$ is reported as a routing diagnostic with the
discrepancy named, and no claim rests on it
(Appendix~\ref{app:tuning-protocol}).

\end{itemize}

This register closes the discussion: the errors the apparatus caught
(Section~\ref{sec:disc-selfcheck}) and the debt it still carries are
two views of the same measurement discipline. What remains is to state
what survives both --- what was demonstrated, on what scope, and what
is handed to future work --- and that is the Conclusion.

%%% The Conclusion and Future Work sections of this chapter continue
%%% in 03_Content/8_Conclusion.tex, which is part of this chapter and
%%% carries no chapter command of its own.
